\documentclass[12pt,a4paper]{article}

\usepackage[utf8]{inputenc}
\usepackage[T1]{fontenc}
\usepackage{lmodern}
\usepackage[english]{babel}
\usepackage[margin=2.5cm]{geometry}
\usepackage{microtype}
\usepackage{setspace}
\usepackage{amsmath}
\usepackage[hidelinks]{hyperref}

\onehalfspacing
\setlength{\parindent}{0pt}

\begin{document}

\section*{Abstract}

Behavioural test batteries obtain several domains of information from the same animals, but an assay placed inside a sequence may measure something subtly different from the same assay performed first. This measurement-context problem is well described in rodents and remains less extensively characterised in zebrafish, particularly for batteries that extend beyond anxiety-related assays. We evaluated a sequential battery in adult zebrafish (\textit{Danio rerio}) comprising the light--dark test, the novel tank test, and a swimming-resistance assay, and compared each measurement against a context in which the assay occupied a different position: a reversed behavioural context, in which the novel tank preceded the light--dark test, and an isolated endurance context. Ethanol (0, 0.5, 1\%) applied for 1, 24, or 96~h served as a secondary multidomain perturbation. Context differences were estimated within matched concentration-by-duration cells using rank-based tests, Hedges' $g$ with bootstrap confidence intervals, and Benjamini--Hochberg adjustment across 117 endpoint-by-cell comparisons; equivalence testing was not conducted because endpoint-specific margins were not prespecified. Ninety-eight of the 117 comparisons (83.8\%) showed no detected adjusted context difference, including 73 of 81 light--dark and novel tank comparisons. Differences were detected in 19 comparisons and were concentrated in the endurance domain (11 of 36), where they were large (median absolute $g$ = 1.16) and ran opposite to a fatigue account: resistance performance was lower in the isolated context than in the sequential battery, while reinsertion attempts were higher. Exploratory correlation and principal component analyses of the same fish showed strong within-domain and weaker cross-domain relationships, with several endpoints carrying overlapping information. Ethanol produced adjusted responses in ten of fourteen modelled endpoints across all assay domains. Measurement-context interference was therefore selective rather than pervasive, supporting continued use of the battery while identifying swimming resistance as the priority for methodological characterisation.

\vspace{1em}
\noindent\textbf{Keywords:} \textit{Danio rerio}; behavioural test battery; measurement context; swimming resistance; ethanol; multidomain phenotyping

\end{document}
